\clearpage
\fancyhead[L]{\small TABLES DE LA LUNE.}
\SourceLeaf{645}
\begin{center}\small Göttingische gelehrte Anzeigen. Stück 23. Seiten 225..227. 7. Februar 1825.\\[3pt]\rule{.72\linewidth}{.3pt}\end{center}

\emph{Tables de la lune, formées par la seule théorie de l'attraction, et suivant la division de la circonférence en \textup{400} degrés. Par le Baron \textup{\textsc{de Damoiseau}}, Lieutenant Colonel etc.} Paris 1824. Bei Bachelier 90 Seiten in 4.

So vollkommen die Mondtafeln durch die Arbeiten von \textsc{Tobias Mayer}, \textsc{Bürg} und \textsc{Burckhardt} auch geworden sind, so blieb ihr empirischer Ursprung doch gewissermaßen ein Vorwurf für die theoretische Astronomie. Die Genauigkeit, mit der man die Ungleichheiten der Mondbewegung bisher aus dem Prinzip der allgemeinen Schwere ableiten konnte, blieb noch deutlich hinter der Genauigkeit zurück, die man auf empirischem Weg erreicht hatte. Hier war noch ein ehrenvoller Kranz zu erringen, indem man sehr verwickelte analytische Rechnungen geschickt anging und sie sorgfältig und unermüdlich durchführte. Die Pariser Akademie ermunterte dazu, indem sie diese Rechnungen zum Gegenstand einer Preisaufgabe machte. Dadurch veranlasste sie die erfolgreichen Arbeiten von \textsc{Damoiseau}, \textsc{Carlini} und \textsc{Plana}. Hoffentlich werden wir bald die theoretischen Untersuchungen dieser verdienten Astronomen erhalten. Zunächst liegen uns hier die Tafeln vor, die auf den Entwicklungen von \textsc{Damoiseau} beruhen und deren Erfolg bereits bezeugen.

Die Einrichtung dieser Tafeln weicht von der bisher gebräuchlichen darin ab, dass sich für die Länge des Mondes in der Ekliptik und für die Horizontalparallaxe sämtliche Argumente proportional zur Zeit ändern, während für die Breite die bereits um die Ungleichheiten berichtigte Länge zugrunde liegt. Diese Einrichtung vermindert zwar die Konvergenz beträchtlich und erfordert eine größere Zahl von Gleichungen (für die Länge sind sie in 46 Tafeln untergebracht, mit dem Beitrag zweier Tafeln für die Störungen durch Venus und Jupiter). Der völlig gleichförmige Ablauf aller Rechnungen gleicht dies nun wohl weitgehend aus, obwohl es bei der Benutzung immer noch etwas unbequem ist, dass viele kleine Gleichungen in Tafeln mit doppelten Eingängen zusammengefasst sind. Hätte man Tafeln mit doppelten Eingängen ganz vermeiden wollen, so wäre die Gesamtzahl der Tafeln für die Länge auf 67 gestiegen. Die Übernahme der von \textsc{Carlini} in seinen Sonnentafeln gewählten Einrichtungen, die tägliche Änderung der gleichförmig wachsenden Argumente als deren Einheit anzusetzen, würde den Gebrauch der Tafeln wohl sehr erleichtern.

Um die mit der Mondtheorie von \textsc{Damoiseau} erreichte Genauigkeit beurteilen und mit der Genauigkeit der früheren Tafeln vergleichen zu können, ist am Schluss des Werks ein Vergleich von 50 Beobachtungen, die während eines anderthalbjährigen Zeitraums 1802 und 1803 an der Sternwarte in Greenwich angestellt wurden, mit den Tafeln von \textsc{Bürg}, \textsc{Burckhardt} und \textsc{Damoiseau} beigefügt. Schon ein Blick zeigt, dass die letzten Tafeln die Beobachtungen nicht schlechter darstellen als die früheren. Um jedoch ein bestimmteres Urteil fällen zu können, hat sich der Rezensent die Mühe gemacht, daraus den sogenannten mittleren Fehler abzuleiten, der sich
\begin{center}
\begin{tabular}{l|c|c}
bei den Tafeln von & in der Länge & in der Breite \\
\hline
Bürg & $5''80$ & $9''18$ \\
Burckhardt & $5.37$ & $7.82$ \\
Damoiseau & $5.23$ & $7.12$ \\
\end{tabular}
\end{center}
\noindent ergeben hat. Soweit sich also aus dieser verhältnismäßig noch kleinen Zahl von Vergleichen schließen lässt, hat die Theorie den Empirismus nun tatsächlich überholt.
Ein nicht ganz \SourceLeaf{646} unbeträchtlicher Teil der Abweichungen mag übrigens wohl den Beobachtungen selbst zuzuschreiben sein.
Nachdem die Bahn aber einmal so weit gebrochen ist, scheint die Hoffnung begründet, dass die Vollkommenheit der Tafeln durch unermüdliches Fortschreiten auf demselben Weg immer weiter zunehmen wird und die Mondbewegung fast mit derselben Schärfe dargestellt werden wird, zu der die verfeinerte Beobachtungskunst fähig ist.
